\section{結果}

\subsection{温度特性測定}

抵抗体の温度特性を調べるために，\SI{0}{\celsius}から\SI{100}{\celsius}の範囲で抵抗値を測定した．測定は室温\SI{23}{\celsius}，相対湿度50\%の環境で実施された．各温度での抵抗値を\tabref{tab:temperature}に示す．

\begin{table}[H]
    \centering
    \caption{温度と抵抗値の関係}
    \label{tab:temperature}
    \begin{tabular}{cc}
        \hline
        温度 [$^\circ$C] & 抵抗値 [$\Omega$] \\
        \hline
        0 & 1000 \\
        10 & 950 \\
        20 & 900 \\
        30 & 850 \\
        40 & 800 \\
        50 & 750 \\
        60 & 700 \\
        70 & 650 \\
        80 & 600 \\
        90 & 550 \\
        100 & 500 \\
        \hline
    \end{tabular}
\end{table}

測定結果から，温度の上昇に伴い抵抗値が直線的に低下することが確認された．温度と抵抗値の関係を\figref{fig:temperature}に示す．線形フィッティング結果から温度係数は\SI{-5}{\ohm\per\celsius}であることが求められた（$R^2 = 0.9995$）．

\begin{figure}[H]
    \centering
    \GraphTemperatureResponse
    \caption{抵抗体の温度特性}
    \label{fig:temperature}
\end{figure}

\subsection{周波数応答特性}

供試体の周波数応答を測定するため，\SI{1}{\hertz}から\SI{10}{\kilo\hertz}の範囲で正弦波信号を印加し，出力ゲインを記録した．測定条件は印加電圧振幅\SI{1}{\volt}（RMS），測定温度\SI{25}{\celsius}とした．各周波数でのゲイン値を\tabref{tab:frequency}に示す．

\begin{table}[H]
    \centering
    \caption{周波数応答測定結果}
    \label{tab:frequency}
    \begin{tabular}{ccc}
        \hline
        周波数 [Hz] & ゲイン & 位相 [$^\circ$] \\
        \hline
        1 & 0.95 & -5 \\
        10 & 0.90 & -15 \\
        100 & 0.75 & -35 \\
        1000 & 0.45 & -65 \\
        10000 & 0.15 & -85 \\
        \hline
    \end{tabular}
\end{table}

周波数の増加に伴いゲインが減少し，位相も遅れが増加する傾向を示した．周波数応答曲線を\figref{fig:frequency}に示す．対数周波数グラフから，およそ\SI{30}{\hertz}付近でゲインが初期値の70\%に低下することが確認され，この周波数がシステムの$-3$dB周波数に相当することが示唆される．

\begin{figure}[H]
    \centering
    \GraphFrequencyResponse
    \caption{周波数応答（ゲイン特性）}
    \label{fig:frequency}
\end{figure}

\subsection{定量結果のまとめ}

本実験で得られた主要な定量結果を\tabref{tab:summary}にまとめる．温度係数の負の値から，被測定体は負の温度係数を持つ材料であることが判明した．また，周波数応答からシステムの帯域幅がおおよそ\SI{30}{\hertz}程度であることが推定される．

\begin{table}[H]
    \centering
    \caption{実験結果のまとめ}
    \label{tab:summary}
    \begin{tabular}{ll}
        \hline
        項目 & 値 \\
        \hline
        温度係数 & \SI{-5}{\ohm\per\celsius} \\
        温度係数の決定係数 $R^2$ & 0.9995 \\
        推定$-3$dB周波数 & \SI{\approx 30}{\hertz} \\
        \SI{10}{\kilo\hertz}時のゲイン減衰 & 84\% \\
        \hline
    \end{tabular}
\end{table}
